% Securities Lending and Portfolio Decisions

We begin our analysis by examining whether mutual funds adjust their portfolio positions in response to securities lending activity. Our baseline specification, described formally in Section~\ref{sec:empirical-strategy}, regresses monthly position changes on lagged lending indicators while absorbing unobserved heterogeneity through a saturated set of fixed effects.

\subsection{Baseline Results}

Table~\ref{table_baseline} presents the baseline results. The dependent variable is $PosChange_{i,j,t}$, the monthly percentage change in the number of shares of stock $i$ held by fund $j$ at time $t$. Our key explanatory variable is $OnLoan_{i,j,t-1}$ (standardized), which measures the dollar value of shares lent by fund $j$ in stock $i$ as a fraction of market capitalization at $t-1$. All specifications include security$\times$time, fund$\times$time, and fund$\times$security fixed effects. Standard errors are clustered at the fund level.

In column~(1), we estimate the relationship for the full sample of funds. The coefficient on $OnLoan_{i,j,t-1}$ is $-0.2069$ and statistically significant at the 1\% level ($t = -4.23$). Funds reduce their positions in stocks they lend. This effect is economically meaningful: a one-standard-deviation increase in lending activity is associated with a 0.21 percentage point reduction in position size in the subsequent month.

\subsection{Active vs.\ Passive Funds}

A core prediction of the information acquisition hypothesis is that only active fund managers---who can freely adjust their portfolios---should respond to the lending signal. Passive fund managers, constrained to track their benchmark index, should not trade on the information even if they observe it.

Column~(2) of Table~\ref{table_baseline} tests this prediction by interacting $OnLoan_{i,j,t-1}$ with an indicator for actively managed funds, as in equation~(\ref{eq:baseline_active}). The interaction term $Active_j \times OnLoan_{i,j,t-1}$ is $-0.3382$ and highly significant ($t = -5.33$), while the uninteracted $OnLoan_{i,j,t-1}$ coefficient becomes small and statistically insignificant ($0.0133$, $t = 0.80$). The position reduction is entirely driven by active funds: passive funds do not adjust their portfolios in response to lending signals.

Columns~(3)--(5) restrict the sample to active funds only. Column~(3) confirms the baseline effect with a coefficient of $-0.3271$ ($t = -5.39$). Column~(4) replaces the continuous lending measure with the extensive margin indicator $D(\mathit{Lender})_{i,j,t-1}$, equal to one if the fund lent any shares of the stock in the prior month. The coefficient of $-1.9165$ ($t = -4.18$) indicates that active lending funds reduce positions significantly more than non-lending active funds. Column~(5) includes both measures simultaneously; the continuous $OnLoan$ measure dominates ($-0.3206$, $t = -5.49$), while the extensive margin indicator becomes insignificant ($-0.0584$, $t = -0.15$), suggesting that the \textit{intensity} of lending---not merely lending participation---drives the trading response.

Taken together, the results in Table~\ref{table_baseline} establish that active fund managers reduce positions in stocks they lend, with the effect concentrated at the intensive margin of lending. The passive fund placebo (column~2) rules out pure mechanical explanations such as share recall for voting purposes, which would affect passive and active funds similarly. This active-passive asymmetry is the central empirical prediction of \citet{Nurisso2026}: active lenders face a commitment problem that arises precisely from their ability to trade on the private signal, leading short sellers to prefer passive lenders as counterparties. The corollary for our tests is clean: passive funds, who cannot exploit the signal, exhibit zero response, while active funds respond strongly.
